\documentclass[10pt,a4paper]{amsart}
\usepackage[margin=19mm]{geometry}
\usepackage{amsmath,amssymb,mathtools}
\usepackage[scr=rsfs,cal=euler]{mathalfa}
\usepackage{xfrac}
\usepackage[unicode,hidelinks,bookmarks=false]{hyperref}
\newcommand{\ie}{i.e.\ }
\newcommand{\cf}{cf.\ }
\newcommand{\resp}{respectively\ }
\newcommand{\Subsection}{\subsection}
\providecommand{\nobreakdash}{\nobreak\hbox{-}}
\setlength{\emergencystretch}{2em}
\multlinegap0pt
\usepackage{amssymb}
\usepackage[shortlabels]{enumitem}
\usepackage{bm}
\newtheorem{prop}{Proposition}
\newtheorem{thm}{Theorem}
\newcommand{\Z}{\mathbb{Z}}
\title{Embeddability between $3$-manifolds is not stable}
\author{Giulio Belletti, Renaud Detcherry}
\date{Source: arXiv:2604.22387v2 (CC BY 4.0)}
\begin{document}
\maketitle

\noindent\emph{Source and licence.} Embeddability between $3$-manifolds is not stable. Version arXiv:2604.22387v2 (CC BY 4.0), distributed by arXiv under the Creative Commons Attribution 4.0 International licence, http://creativecommons.org/licenses/by/4.0/, as recorded in the arXiv licence metadata for this exact version and shown on the abstract page https://arxiv.org/abs/2604.22387v2. Full licence notice: \texttt{licenses/arxiv-2604.22387-CC-BY-4.0.txt}.\par\smallskip

\noindent\textit{Editorial scope.} Counted unit: the article's thm thm:main (source0.tex lines 588--597). The article's complete original proof of it is delivered with it (source0.tex lines 627--646, one \texttt{proof} environment of 20 lines, which the article places away from the statement). No mathematics is written, repaired or re-derived: the statement and the proof are the article's own lines, in the article's own notation, and no retained line is substituted or re-flowed.  Retained uncounted as the source-local context the proof needs: lines 313--324.  Nothing was omitted from any retained span, so the package carries no omission record.  No retained line contains a citation, so the wrapper carries no bibliography and no bibliography entry.  No retained line contains a figure environment, an image-inclusion command or drawing code, so no image is delivered and the package declares no asset dependency.  Editorial formatting only, in addition: an AMS \texttt{amsart} preamble that restates the article's own declarations and macros used by the retained lines, namely the environments \texttt{prop}, \texttt{thm} on separate counters, together with the standard packages those lines need.  The retained environments print in this document as Proposition 1, Theorem 1 in order of appearance, whereas the article prints the same items with the numbering of its own full document.  The complete native source of the article as distributed in the arXiv e-print for this exact version is bundled unchanged as \texttt{sources/199087/source0.tex}.
\begin{prop}\label{prop:VeryGood}
	A compact manifold $M$ is very good if any of the following conditions holds:
	\begin{enumerate}
		\item $M$ is a rational homology sphere;
		\item $M=M_1\# M_2$ and both $M_1,M_2$ are very good;
		\item $M$ is the double of a very good manifold;
		\item $M$ is a submanifold of a very good manifold; 
		\item $M$ is obtained from a very good manifold by attaching a connected Seifert cobordism with one input and at least one output boundary components.
	\end{enumerate}
	
	Furthermore, if $M$ is hyperbolic, for any prime $p,$ there exists a finite cover of $M$ which is $p$-good.
\end{prop}

\begin{thm}\label{thm:main}
	Take $k\geq 1$, $p\geq 5$ prime, $n$ coprime to $p$, $\mathcal{G}=\{G_1,\dots, G_r\}$ a finite collection of finite groups, $M$ a compact oriented $p$-good $3$-manifold, $g$ a positive integer and $v>0$. Let $N$ be any compact oriented manifold; then there is a hyperbolic $3$-manifold $N'$ with the following properties:
	
	\begin{itemize}
		\item  $N'$ is $(Y_k,T_n,\mathcal{G})$-equivalent to $N$;
		\item $N'$ does not contain any essential, non-boundary parallel surface of genus less than $g$;
		\item $N'$ has volume larger than $v$;
		\item $I_p(M)\nsubseteq I_p(N')$ which implies that $N'$ does not embed in $M$.
	\end{itemize}
\end{thm}

\begin{proof}
	(1) Let $p_1,\dots, p_l$ be primes that do not divide $n$ and such that $M$ is $p_i$-good, and call their product $d$.  By Theorem \ref{thm:main}, for each $i$ there is a hyperbolic $M_i$ which is $(Y_k,T_{nd/p_i})$-equivalent to $M$ but such that $I_{p_i}(M)\nsubseteq I_{p_i}(M_i)$. However, since $M_i$ is $T_{p_j}$-equivalent to $M$, $I_{p_j}(M)=I_{p_j}(M_i)$ for all $i\neq j$, which implies that $I_{p_i}(M_i)\nsubseteq I_{p_i}(M_j)$ and therefore $M_i$ does not embed in $M_j$.
	
	(2) We first note that any $3$-manifold which is $Y_1$-equivalent to a subspace of rational homology sphere is itself a subspace of a rational homology sphere. We construct the sequence $(M_i)_{i\geq 1},$ as well as a sequence of prime numbers $(p_i)_{i\geq 1}$ coprime to $n$ inductively, so that 
	\[\forall j<i, \ I_{p_j}(M_i)=I_{p_j}(M)\]
	and
	\[ I_{p_i}(M_1)I_{p_i}(M_2)\ldots I_{p_i}(M_{i-1})I_{p_i}(M) \nsubseteq I_{p_i}(M_i).\]
	Note that these conditions immediately imply that $M_i$ does not embed in $M_j$ for $i\neq j.$
	
	
	To start with, since $M$ is very good by Proposition \ref{prop:VeryGood}, let $p_1\geq 5$ be a prime such that $M$ is $p_1$-good. By Theorem \ref{thm:main}, there exists a hyperbolic $M_1$ which is $(Y_k,T_n)$-equivalent to $M,$ such that $I_{p_1}(M)\nsubseteq I_{p_1}(M_1).$
	
	 Assume that $M_1,\ldots, M_l$ have been constructed. Since they are all $(Y_k,T_n)$-equivalent to $M,$ they are all $Y_1$-equivalent to $M,$ hence they are all subspaces of rational homology spheres, and by Proposition \ref{prop:VeryGood}, they are all very good.
	
	Let $p_{l+1}\geq 5$ be a prime which does not divide $N=np_1\ldots p_l,$ and such that $RT_{p_{l+1}}(M)\neq 0$ and that $RT_{p_{l+1}}(M_i)\neq 0$ for any $i\leq l.$ Then the ideal 
	$$J=I_{p_{l+1}}(M_1)I_{p_{l+1}}(M_2)\ldots I_{p_{l+1}}(M_{l})I_{p_{l+1}}(M)$$ of $\Z[\zeta_{p_{l+1}},\frac{1}{p_{l+1}}]$ is non-zero. By the proof of Theorem \ref{thm:main}, there exists a hyperbolic $3$-manifold $M_{l+1}$ that is $(Y_k,T_N)$-equivalent to $M,$ and such that
	$$J\nsubseteq I_{p_{l+1}}(M_{l+1}).$$
	
	This completes the induction step; we have constructed an infinite sequence of $3$-manifolds that are all $(Y_k,T_n)$-equivalent to $M$ but pairwise do not embed in each other.
\end{proof}

\end{document}
